\documentclass[12pt]{article}

% ============================================================
%  Core math
% ============================================================
\usepackage{amsmath, amssymb, amsthm}
\usepackage{mathtools}
\usepackage{array}
\usepackage{longtable}
% Ragged-right fixed-width column (avoids underfull-hbox warnings in p-columns)
\newcolumntype{P}[1]{>{\raggedright\arraybackslash}p{#1}}

% Diagrams
\usepackage{tikz}
\usepackage{tikz-cd}
\usetikzlibrary{arrows.meta, positioning, calc}

% References
\usepackage{hyperref}
\usepackage{cleveref}

% Graphics
\usepackage{graphicx}

% Page layout
\usepackage[margin=1in]{geometry}

% GrokRxiv sidebar
\usepackage{everypage}
\usepackage{xcolor}

\hypersetup{
  colorlinks=true,
  linkcolor=blue!55!black,
  citecolor=green!45!black,
  urlcolor=blue!55!black
}

% ============================================================
%  Theorem environments
% ============================================================
\newtheorem{theorem}{Theorem}[section]
\newtheorem{proposition}[theorem]{Proposition}
\newtheorem{lemma}[theorem]{Lemma}
\newtheorem{corollary}[theorem]{Corollary}
\theoremstyle{definition}
\newtheorem{definition}[theorem]{Definition}
\newtheorem{example}[theorem]{Example}
\newtheorem{axiom}[theorem]{Axiom}
\theoremstyle{remark}
\newtheorem{remark}[theorem]{Remark}

% ============================================================
%  Macros
% ============================================================
\newcommand{\Spec}{\operatorname{Spec}}
\newcommand{\Proj}{\operatorname{Proj}}
\newcommand{\Mot}{\operatorname{Mot}}
\newcommand{\Real}{\operatorname{Real}}
\newcommand{\Obs}{\operatorname{Obs}}
\newcommand{\Dom}{\mathsf{Dom}}
\newcommand{\Math}{\mathsf{Math}}
\newcommand{\Phys}{\mathsf{Phys}}
\newcommand{\Sch}{\mathsf{Sch}}
\newcommand{\Set}{\mathsf{Set}}
\newcommand{\Grpd}{\mathsf{Gpd}}
\newcommand{\Aff}{\mathsf{Aff}}
\newcommand{\Ring}{\mathsf{Ring}}
\newcommand{\CRing}{\mathsf{CRing}}
\newcommand{\Vect}{\mathsf{Vect}}
\newcommand{\Ideal}{\mathfrak{a}}
\newcommand{\Crit}{\operatorname{Crit}}
\newcommand{\dCrit}{\mathbf{dCrit}}
\newcommand{\Stab}{\operatorname{Stab}}
\newcommand{\Aut}{\operatorname{Aut}}
\newcommand{\Res}{\operatorname{Res}}
\newcommand{\HH}{\mathbb{H}}
\newcommand{\CC}{\mathbb{C}}
\newcommand{\RR}{\mathbb{R}}
\newcommand{\ZZ}{\mathbb{Z}}
\newcommand{\QQ}{\mathbb{Q}}
\newcommand{\PP}{\mathbb{P}}
\newcommand{\AAf}{\mathbb{A}}
\newcommand{\Ohm}{\mathcal{O}}
\newcommand{\per}{\operatorname{per}}
\newcommand{\Pic}{\operatorname{Pic}}
\newcommand{\Ext}{\operatorname{Ext}}
\newcommand{\Hom}{\operatorname{Hom}}
\newcommand{\rk}{\operatorname{rank}}
\newcommand{\amp}{\mathcal{A}}
\newcommand{\coeff}{\mathbb{k}}
\newcommand{\op}{\mathrm{op}}
\newcommand{\GM}{\nabla_{\mathrm{GM}}}
\newcommand{\GMd}{\nabla^{\mathrm{GM}}_{\partial}}
\newcommand{\Fil}{F^{\bullet}}
\newcommand{\status}[1]{\textnormal{\textsf{\textbf{[#1]}}}}

% GrokRxiv sidebar
\definecolor{grokgray}{RGB}{110,110,110}
\AddEverypageHook{%
  \ifnum\value{page}=1
    \begin{tikzpicture}[remember picture, overlay]
      \node[
        rotate=90,
        anchor=south,
        font=\Large\sffamily\bfseries\color{grokgray},
        inner sep=0pt
      ] at ([xshift=38pt, yshift=0.52\paperheight]current page.south west)
      {GrokRxiv:2026.07.algebraic-geometry-physical-possibility\quad
       [\,math.AG\,]\quad
       04 Jul 2026};
    \end{tikzpicture}
  \fi
}

% ============================================================
\title{\textbf{Algebraic Geometry: Spaces of Physical Possibility}\\[4pt]
\large Part III of \emph{A Math$\to$Physics Representation Library}}

\author{Matthew Long \\
\textit{The YonedaAI Collaboration} \\
\textit{YonedaAI Research Collective} \\
Chicago, IL \\
\texttt{research@yonedaai.com} $\cdot$ \url{https://yonedaai.com}}

\date{4 July 2026}

\begin{document}
\maketitle

\begin{abstract}
This is Part III of the modular series \emph{A Math$\to$Physics Representation Library}, which
formalizes each mathematical domain as a physical-representation module and composes the modules
hierarchically. Part~I introduced the \emph{representation stack} and the realization pipeline
$\Obs_\alpha(M) = \Obs(\Real_\alpha(\Phi(M)))$; Part~II gave observables an internal
\emph{coalgebraic anatomy} through motives, periods, and the Goncharov coproduct. Here we supply
the geometric \emph{stage} on which those motives and periods live and vary: the theory of schemes,
moduli spaces and stacks, variations of Hodge structure, and positive geometries. Our organizing
thesis is that \emph{algebraic geometry is the geometry of physical possibility}: an affine variety
is a classical solution space, a scheme is a solution space enriched with infinitesimal structure,
a moduli space is a space of physically distinct configurations, and a moduli \emph{stack} is such a
space with its gauge automorphisms remembered rather than quotiented away. We work throughout with the
epistemic status discipline of Part~I, labelling every dictionary entry \status{S} (standard),
\status{H} (heuristic), or \status{P} (speculative). We prove four theorems that upgrade Part~II's
per-object realization pipeline to a family-indexed one: (T1) Gauss--Manin flat transport realizes the
Part~II period functor over an entire moduli family, with amplitudes constrained by the associated
Picard--Fuchs system; (T2) a moduli stack $[X/G]$ is a strictly finer representation object than its
coarse space exactly when stabilizers are nontrivial, giving a rigorous algebro-geometric form of
Part~I's ``stacks retain gauge information'' principle; (T3) the recursive residue structure of a
positive geometry reproduces a coalgebra-like decomposition tree conjecturally isomorphic to Part~II's
motivic coaction tree (flagged \status{H}); and (T4) the derived critical locus formalizes the on-shell
observable algebra $R/I$ at the Batalin--Vilkovisky level, upgrading a heuristic dictionary entry to a
status-\status{S} statement via shifted symplectic geometry. We accompany the paper with executable
\textsf{Haskell} encodings of schemes-as-functors-of-points, quotient stacks, variations of Hodge
structure, and positive geometries, together with a best-effort \textsf{Lean~4} sketch of the same
type signatures. Part~III thereby sets the geometric foundation on which Part~IV (algebraic topology:
conserved global information) will build.
\end{abstract}

\tableofcontents

% ============================================================
\section{Introduction}
\label{sec:intro}

\subsection{The series and its modular ladder}

The present paper is the third module of a seven-paper program, \emph{A Math$\to$Physics
Representation Library}. The program takes as its point of departure the observation that many
physical quantities are not merely \emph{described by} mathematics but are literally \emph{obtained
as realizations} of structured mathematical information~\cite{primary}. Rather than assert one
universal functor from all of mathematics to all of physics (an object we argue in
\Cref{sec:limitations} is unlikely to be well defined), the program is deliberately
\emph{modular}: each mathematical domain is formalized as its own physical-representation module,
each module is defensible on its own terms, and the modules compose hierarchically so that new
structure and new emergent physical content appear at each level.

The ladder proceeds as follows.
\begin{itemize}
  \item \textbf{Part~I (Foundations).} Introduces the category $\Dom$ of mathematical domains, the
  categories $\Math_d$ of structures in each domain, the target $\Phys$ of physical representations,
  the notion of a \emph{representation entry} $E=(M,P,\tau,\sigma)$, the \emph{representation
  prestack} and \emph{representation stack}, the four foundational axioms (Realization,
  Equivalence, Locality/descent, Decomposition), and the epistemic status system \status{S}/\status{H}/\status{P}.
  Its central formula is the \emph{realization pipeline}
  \begin{equation}\label{eq:pipeline}
    \Obs_\alpha(M) \;=\; \Obs\bigl(\Real_\alpha(\Phi(M))\bigr).
  \end{equation}
  \item \textbf{Part~II (Motives, periods, amplitudes).} Makes the abstract realization channels
  $\Real_\alpha$ of Part~I concrete: Betti, de Rham, Hodge, and \'etale realizations of motives; the
  period datum $\Pi=(X,D,\omega,\Gamma)$ with $\per(\Pi)=\int_\Gamma\omega$; the motivic amplitude
  object $\amp^{\mathfrak m}$ with $\per(\amp^{\mathfrak m})=\amp$; and the coalgebraic anatomy
  encoded by the coproduct $\Delta$ and cobracket $\delta$. Observables acquire an internal
  \emph{decomposability} into primitive informational pieces.
  \item \textbf{Part~III (this paper).} Supplies the geometric \emph{stage}. Part~II's motives and
  periods do not live in a vacuum: a period datum requires $X$ to be an actual algebraic or analytic
  space, and a \emph{family} of amplitudes over kinematic space is a variation of Hodge structure over
  a moduli space. We formalize schemes, moduli spaces and stacks, Hodge theory and Gauss--Manin
  transport, and positive geometries, and we show that Part~II's per-object pipeline extends to a
  \emph{family} pipeline fibered over the geometry.
  \item \textbf{Parts~IV--VII.} Algebraic topology (conserved global information), category
  theory and homotopy type theory (composition and identity), sheaves/stacks/gauge redundancy/quantum
  high availability, and finally a synthesis paper recomposing the ladder. Part~III sets up Part~IV by
  producing the varieties and moduli whose topological invariants Part~IV will study.
\end{itemize}

\subsection{Thesis of Part~III}

Our thesis, following the primary source's algebraic-geometry appendix, is compact:
\begin{quote}
\emph{Algebraic geometry is the geometry of physical possibility.}
\end{quote}
The word ``possibility'' is meant in a precise, layered sense that the machinery of algebraic
geometry makes available:
\begin{enumerate}
  \item \textbf{Solutions as points.} The set of states satisfying a system of polynomial constraints
  is an \emph{affine variety} $V(\Ideal)$; its coordinate ring $R/\Ideal$ is the algebra of on-shell
  observables. \status{S}\ as algebra; the reading ``measurements generate the space of
  possibilities'' is \status{H}.
  \item \textbf{Infinitesimal possibility.} Passing from varieties to \emph{schemes} $\Spec R$ retains
  nilpotent (infinitesimal, off-shell, ghost-like) directions that varieties discard; possibility
  acquires a tangent/deformation structure.
  \item \textbf{Families of possibility.} A \emph{moduli space} parametrizes physically distinct
  configurations; the space itself, its singularities, and its compactification become physical
  objects (vacua, thresholds, asymptotic sectors).
  \item \textbf{Possibility modulo redundancy.} A \emph{moduli stack} $[X/G]$ remembers the gauge
  automorphisms (stabilizers) of each configuration, distinguishing genuine physical states from
  redundant descriptions, which is the geometric incarnation of Part~I's Equivalence axiom.
  \item \textbf{Possibility varying analytically.} A \emph{variation of Hodge structure} equips a
  family with a flat Gauss--Manin connection whose flat sections are the Betti cycles; pairing them
  against the (non-flat) holomorphic form gives the amplitudes across kinematic space, which obey the
  Picard--Fuchs differential equations.
  \item \textbf{Possibility with a boundary.} A \emph{positive geometry} carries a canonical form whose
  boundary residues recursively encode lower-dimensional geometries, a geometric model of physical
  factorization and singularity structure.
\end{enumerate}

\subsection{Contributions}

\begin{enumerate}
  \item A self-contained formalization (\Cref{sec:framework,sec:schemes}) of the algebraic-geometry
  layer of the representation library: schemes via the functor of points, ideals as constraints, the
  on-shell algebra $R/\Ideal$, and the full \status{S}/\status{H}/\status{P}-labelled dictionary of
  \Cref{sec:dictionary}, reproduced and organized from the primary source.
  \item Four theorems with proofs (\Cref{sec:theorems}), each of which extends a Part~II construction
  over a geometric family: Gauss--Manin family realization (\Cref{thm:gm}), stack-vs-coarse strictness
  (\Cref{thm:stack}), positive-geometry residue/coaction correspondence (\Cref{thm:posgeom},
  status \status{H}), and the derived critical locus as BV on-shell algebra (\Cref{thm:dcrit}).
  \item An explicit family realization pipeline (\Cref{sec:family}) generalizing
  \eqref{eq:pipeline} from a single mathematical object to a moduli-indexed family, together with a
  status-propagation calculus inherited from Part~I.
  \item Executable \textsf{Haskell} encodings and a best-effort \textsf{Lean~4} sketch
  (\Cref{sec:code}) of schemes-as-functors-of-points, quotient stacks with stabilizer data,
  variations of Hodge structure with Gauss--Manin connections, and positive geometries with recursive
  residue trees.
  \item An explicit accounting (\Cref{sec:discussion}) of how Part~III respects the program's five
  standing limitations, and a list of open problems, notably the conjectural identification of
  positive-geometry residue trees with motivic coaction trees, handed forward to Parts~IV--VII.
\end{enumerate}

\subsection{Epistemic status discipline}
\label{sec:status-intro}

We preserve verbatim the three-level status system of Part~I.

\begin{center}
\begin{tabular}{c P{5.6cm} P{5.4cm}}
\hline
\textbf{Label} & \textbf{Meaning} & \textbf{Canonical example} \\
\hline
\status{S} & standard use in mathematics or mathematical physics & TQFT as a functor $\mathsf{Bord}_n\to\Vect$ \\
\status{H} & strong heuristic representation dictionary entry & motive as ``functional essence'' \\
\status{P} & speculative philosophical ontology & reality as motivic-categorical realization \\
\hline
\end{tabular}
\end{center}

Composite labels (\status{S/H}, \status{H/P}) appear where an entry is standard as pure mathematics
but heuristic or speculative in its \emph{physical} reading; we preserve these composite labels
verbatim. This discipline is not decoration: it is the program's central epistemic-honesty device, and
in \Cref{sec:theorems} we state precisely which of our four theorems are status-\status{S}
(established mathematics applied to a physical reading), which are status-\status{H} (heuristically
plausible, original to this program), and which are status-\status{P}.

% ============================================================
\section{Mathematical Framework}
\label{sec:framework}

We recall the ambient framework of Parts~I--II in the compressed form we need, then specialize to the
algebraic-geometry domain $d = \mathsf{AG}$.

\subsection{Domains, structures, and representation entries}

\begin{definition}[Representation entry, after Part~I]
\label{def:repentry}
Fix a category $\Dom$ of mathematical domains and, for each $d\in\Dom$, a (possibly higher) category
$\Math_d$ of mathematical structures in $d$, together with a category $\Phys$ of physical
representations. A \emph{representation entry} is a quadruple $E=(M,P,\tau,\sigma)$ where
$M\in\Math_d$, $P\in\Phys$, $\tau$ is a semantic translation datum (a functor, a natural
transformation, an equivalence class of models, a physical interpretation map, or a weaker relation),
and $\sigma\in\{\status{S},\status{H},\status{P}\}$ is the epistemic status label.
\end{definition}

For this paper $d=\mathsf{AG}$ and $\Math_{\mathsf{AG}}$ is (a chosen full subcategory of) the category
of schemes, algebraic spaces, and stacks over a fixed base field $\coeff$, together with the auxiliary
categories of Hodge structures and positive geometries introduced below.

\begin{definition}[Status monoid, after Part~I]
\label{def:statusmonoid}
Order the labels by $\status{S} < \status{H} < \status{P}$ and equip $\{\status{S},\status{H},\status{P}\}$
with the binary operation $\sigma_1\wedge\sigma_2 := \max(\sigma_1,\sigma_2)$ (the \emph{less reliable}
of the two). This is a commutative monoid with unit $\status{S}$. For composite labels we take the join
componentwise; e.g.\ $\status{S/H}\wedge\status{S}=\status{S/H}$ and $\status{S/H}\wedge\status{H}=\status{H}$.
\end{definition}

\begin{remark}
\Cref{def:statusmonoid} operationalizes Part~I's status system into a calculus: a chain of translations
is only as reliable as its weakest link, so the status of a composite entry is the $\wedge$-join of the
component statuses. We use this repeatedly to certify that no theorem in \Cref{sec:theorems} silently
launders a heuristic step into a standard conclusion.
\end{remark}

\subsection{The realization pipeline}

Recall the universal formula \eqref{eq:pipeline}. In Part~II the channels $\alpha$ were the
cohomological realizations $\{\text{Betti},\text{de Rham},\text{Hodge},\text{\'etale}\}$ and the abstract
information functor was $\Phi = \Mot$, sending a variety to its motive. The realization pipeline for a
single variety $X$ is the chain
\begin{equation}\label{eq:motchain}
  X \;\xmapsto{\ \Mot\ }\; \Mot(X) \;\xmapsto{\ \Real_\alpha\ }\; \Real_\alpha(\Mot(X)) \;\xmapsto{\ \Obs\ }\; \text{observable}.
\end{equation}
Part~III's central move is to let $X$ \emph{vary}: we replace the single object $X$ by a family
$\pi\colon \mathcal X\to S$ over a base scheme (moduli space) $S$ and ask how \eqref{eq:motchain}
behaves as we move in $S$. The answer, made precise in \Cref{thm:gm}, is that the realization assembles
into a variation of Hodge structure over $S$: the Betti cycles are the flat sections of the Gauss--Manin
connection, and the observable is their pairing with the (non-flat) holomorphic form, a period
function whose variation is governed by $\GM$ and which satisfies the Picard--Fuchs equations.

\subsection{The master diagram, fibered over a base}

\Cref{fig:master} exhibits the family realization pipeline as a commuting diagram. The top row is
Part~I's master diagram for a single object; the bottom row is its family version, and the vertical
arrows are the fibers of the family over a point $s\in S$.

\begin{figure}[!htbp]
\centering
\begin{tikzcd}[column sep=large, row sep=large]
\mathcal X \arrow[r, "\Mot"] \arrow[d, "\pi"'] &
\underline{\Mot}(\mathcal X/S) \arrow[r, "\Real_\alpha"] \arrow[d] &
\mathcal H_\alpha \arrow[r, "\Obs=\int_{\gamma}"] \arrow[d, "\text{fiber}"] &
\text{period } \Pi(s) \arrow[d] \\
S \arrow[r, equal] & S \arrow[r, equal] & S \arrow[r, "\GM"'] & (\text{Picard--Fuchs})
\end{tikzcd}
\caption{The family realization pipeline of Part~III. Over each $s\in S$ the fiber recovers Part~II's
single-object chain \eqref{eq:motchain}; globally the realization $\Real_\alpha$ produces a variation of
Hodge structure $\mathcal H_\alpha$ carrying the Gauss--Manin connection $\GM$. The observable
$\Obs=\int_\gamma$ pairs the (non-flat) holomorphic section of $\mathcal H_\alpha$ against the flat Betti
cycles $\gamma(s)$, yielding the period functions $\Pi(s)$ --- the amplitudes over kinematic space ---
which solve the Picard--Fuchs system (\Cref{thm:gm}). The flat sections of $\GM$ are the Betti cycles,
not the periods.}
\label{fig:master}
\end{figure}

% ============================================================
\section{Schemes as Spaces of Physical Possibility}
\label{sec:schemes}

\subsection{From constraints to varieties}

\begin{definition}[Affine variety, on-shell algebra]
\label{def:variety}
Let $\coeff$ be a field and $R=\coeff[x_1,\dots,x_n]$. For an ideal $\Ideal\subseteq R$ generated by
polynomials $f_1,\dots,f_r$ (the \emph{constraint equations}), the \emph{affine algebraic set} is the
zero locus
\[
  V(\Ideal) \;=\; \{\, a\in \coeff^n : f_i(a)=0,\ i=1,\dots,r \,\},
\]
and the \emph{on-shell observable algebra} is the quotient ring $R/\Ideal$.
\end{definition}

\begin{remark}[Terminology]
In the strict algebro-geometric sense an \emph{affine variety} is an \emph{irreducible} affine algebraic
set, i.e.\ $\Ideal$ is a prime ideal. Throughout this paper we use ``variety'' in the looser, physically
natural sense of an arbitrary solution space $V(\Ideal)$ (an affine algebraic set), since a physical
constraint system need not cut out an irreducible locus. Where irreducibility is actually used
(e.g.\ for a connected smooth family in \Cref{thm:gm}) we say so explicitly.
\end{remark}

Physically, $V(\Ideal)$ is the space of classical states satisfying the imposed equations of motion or
conservation constraints, and $R/\Ideal$ is the algebra of functions (observables) modulo those
constraints: an observable that differs from another by an element of $\Ideal$ takes the same value on
every physical state, hence \emph{is} the same observable on-shell.

\begin{example}[On-shell observable algebra, primary-source Example, \status{S}/\status{H}]
\label{ex:onshell}
Let $R$ be an algebra of fields and coordinates and let $\Ideal$ be the ideal generated by the
equations of motion or constraints. Then $R/\Ideal$ is the natural algebraic representation of
observables modulo the imposed physical equations. \emph{Status:} \status{S}\ as pure algebra; the
ontological reading ``measurements generate the space of physical possibility'' is \status{H}.
\end{example}

\begin{example}[Free particle constraint surface]
Take $R=\RR[q,p,E]$ and $\Ideal=(E-\tfrac{1}{2m}p^2)$. Then $V(\Ideal)\subseteq \RR^3$ is the
paraboloid of on-shell energy-momentum data, and $R/\Ideal\cong \RR[q,p]$: the energy $E$ is not an
independent observable on-shell but a function of $p$. This is the algebraic content of ``imposing the
mass-shell condition.''
\end{example}

\subsection{Schemes and infinitesimal possibility}

Varieties see only reduced structure: they cannot distinguish $V(x)$ from $V(x^2)$, both being the
single point $\{0\}$. Physics, however, frequently needs the infinitesimal or off-shell data that
nilpotents encode. Schemes retain this data.

\begin{definition}[Affine scheme]
\label{def:affsch}
For a commutative ring $R$, the \emph{affine scheme} $\Spec R$ is the set of prime ideals of $R$
equipped with the Zariski topology (closed sets $V(\Ideal)=\{\mathfrak p : \Ideal\subseteq\mathfrak
p\}$) and the structure sheaf $\Ohm_{\Spec R}$ whose stalk at $\mathfrak p$ is the localization
$R_{\mathfrak p}$.
\end{definition}

\begin{definition}[Scheme]
\label{def:scheme}
A \emph{scheme} is a locally ringed space $(X,\Ohm_X)$ that is locally isomorphic to an affine scheme:
every point has an open neighbourhood $U$ with $(U,\Ohm_X|_U)\cong (\Spec R,\Ohm_{\Spec R})$ for some
commutative ring $R$ (see Hartshorne~\cite{hartshorne}, Ch.~II).
\end{definition}

\begin{remark}[Functor of points]
\label{rem:fop}
Equivalently, by Yoneda, a scheme $X$ is determined by its \emph{functor of points}
$h_X\colon \CRing\to\Set$, $R\mapsto \Hom_{\Sch}(\Spec R, X)$; an $R$-point of $X$ is a
$\coeff$-algebra homomorphism from the coordinate ring of $X$ to $R$. This viewpoint is the one we
encode in \textsf{Haskell} and \textsf{Lean} in \Cref{sec:code}: a scheme becomes a rule assigning to
each ring of ``test parameters'' the set of configurations parametrized by that ring. The nilpotent
directions are visible precisely because $R$ is allowed to be nonreduced (e.g.\ dual numbers
$\coeff[\varepsilon]/(\varepsilon^2)$, whose points detect tangent vectors).
\end{remark}

\begin{example}[Nilpotents as virtual directions, \status{H}]
The scheme $\Spec\bigl(\coeff[x]/(x^2)\bigr)$ has one point but a nonzero tangent space: its
$\coeff[\varepsilon]/(\varepsilon^2)$-points are $\{0,\varepsilon\}$, detecting a first-order
deformation invisible to the underlying variety. In the physical dictionary a nilpotent element is an
\emph{infinitesimal thickening} --- a virtual, off-shell, or ghost-like direction (status \status{H}).
This is the algebraic seed of the Batalin--Vilkovisky ghost directions we recover rigorously in
\Cref{thm:dcrit}.
\end{example}

\subsection{Projective varieties and scale invariance}

\begin{definition}[Projective variety]
For a homogeneous ideal $\Ideal\subseteq \coeff[x_0,\dots,x_n]$, the \emph{projective variety}
$V_+(\Ideal)=\Proj\bigl(\coeff[x_0,\dots,x_n]/\Ideal\bigr)\subseteq \PP^n_\coeff$ is the zero locus of
$\Ideal$ in projective space, i.e.\ the space of solutions taken modulo the rescaling
$(x_0,\dots,x_n)\sim (\lambda x_0,\dots,\lambda x_n)$, $\lambda\in\coeff^\times$.
\end{definition}

Physically a projective variety is a \emph{scale-invariant configuration space}: physical data
considered modulo an overall rescaling. Cross-ratios and other conformal invariants live naturally on
projective and Grassmannian geometries; this is the setting in which positive geometries
(\Cref{sec:posgeom}) are formulated.

% ============================================================
\section{The Algebraic Geometry Dictionary}
\label{sec:dictionary}

We reproduce and organize the algebraic-geometry library of the primary source (its Appendix~C),
preserving every status label. The dictionary is the data of a family of representation entries
(\Cref{def:repentry}) in the domain $d=\mathsf{AG}$; the theorems of \Cref{sec:theorems} pick out the
load-bearing rows and prove precise statements about them.

\begin{center}
\small
\begin{longtable}{P{3.6cm} P{4.1cm} P{4.1cm} c}
\caption{Algebraic geometry library (part 1 of 2): schemes, bundles, quotients. Status labels
preserved verbatim from the primary source.}
\label{tab:dict1} \\
\hline
\textbf{Algebraic geometry} & \textbf{Physical representation} & \textbf{Semantic meaning} & \textbf{Status} \\
\hline
\endfirsthead
\multicolumn{4}{c}{\tablename\ \thetable\ -- \textit{continued from previous page}} \\
\hline
\textbf{Algebraic geometry} & \textbf{Physical representation} & \textbf{Semantic meaning} & \textbf{Status} \\
\hline
\endhead
\hline \multicolumn{4}{r}{\textit{continued on next page}} \\
\endfoot
\hline
\endlastfoot
Affine variety & Classical solution space & Allowed states satisfying polynomial constraints & \status{S} \\
Projective variety & Scale-invariant configuration space & Physical data modulo rescaling & \status{S} \\
Scheme & Generalized solution space & Includes infinitesimal/nilpotent structure & \status{S/H} \\
$\Spec R$ & Space encoded by observables & Observables determine geometry & \status{S/H} \\
Coordinate ring & Algebra of observables/functions & Measurements generate space & \status{S} \\
Ideal $\Ideal$ & Constraint equations & Equations of motion/conservation constraints & \status{S/H} \\
Quotient ring $R/\Ideal$ & On-shell observable algebra & Functions modulo physical constraints & \status{S} \\
Nilpotent element & Infinitesimal thickening & Virtual/off-shell/ghost-like direction & \status{H} \\
Sheaf & Locally defined data & Local fields glued into global fields & \status{S} \\
Stalk & Local data at a point & Infinitesimal local physical information & \status{S/H} \\
\v{C}ech cocycle & Gluing rule & Gauge transition function & \status{S} \\
Sheaf cohomology & Obstruction to global gluing & Global anomaly or obstruction & \status{S/H} \\
Line bundle & Phase/gauge degree over space & $U(1)$-type local phase structure & \status{S} \\
Vector bundle & Internal degrees over spacetime & Matter or gauge field carrier & \status{S} \\
Principal $G$-bundle & Gauge configuration & Local symmetry fiber at each point & \status{S} \\
Connection & Gauge potential & Rule for parallel transport & \status{S} \\
Curvature & Field strength & Failure of transport to commute & \status{S} \\
Divisor & Codimension-one locus & Defect, pole, boundary, charge surface & \status{H} \\
Blowup & Resolution of singularity & Regularization/zooming into divergence & \status{S/H} \\
Exceptional divisor & New boundary after resolution & Counterterm or hidden boundary sector & \status{H} \\
Compactification & Add boundary at infinity & Include asymptotic states/IR sectors & \status{S/H} \\
Moduli space & Space of inequivalent structures & Physically distinct vacua/solutions & \status{S} \\
Moduli stack & Moduli with automorphisms retained & Gauge quotient preserving stabilizers & \status{S} \\
Quotient stack $[X/G]$ & Gauge quotient & Physical states modulo symmetry, with isotropy & \status{S} \\
\end{longtable}
\end{center}

\begin{center}
\small
\begin{longtable}{P{3.6cm} P{4.1cm} P{4.1cm} c}
\caption{Algebraic geometry library (part 2 of 2): derived, Hodge-theoretic, and positive-geometric
entries. Status labels preserved verbatim from the primary source.}
\label{tab:dict2} \\
\hline
\textbf{Algebraic geometry} & \textbf{Physical representation} & \textbf{Semantic meaning} & \textbf{Status} \\
\hline
\endfirsthead
\multicolumn{4}{c}{\tablename\ \thetable\ -- \textit{continued from previous page}} \\
\hline
\textbf{Algebraic geometry} & \textbf{Physical representation} & \textbf{Semantic meaning} & \textbf{Status} \\
\hline
\endhead
\hline \multicolumn{4}{r}{\textit{continued on next page}} \\
\endfoot
\hline
\endlastfoot
Derived scheme/stack & Equations plus higher constraints & BV/BRST-style field space with ghosts & \status{S/H} \\
Derived critical locus & Critical points with obstruction data & Perturbative field theory around action & \status{S/H} \\
Tangent complex & Linearized deformation data & Fluctuations around a solution & \status{S} \\
Cotangent complex & Moment/covector deformation data & Phase-space infinitesimals & \status{S/H} \\
Symplectic variety & Geometric phase space & Classical mechanics structure & \status{S} \\
Poisson variety & Bracketed observable space & Classical observable algebra & \status{S} \\
Calabi--Yau variety & Special geometry/compactification & String background or period source & \status{S} \\
Elliptic curve & Torus-like period geometry & Elliptic Feynman integral sector & \status{S} \\
Algebraic curve & Spectral curve & Dispersion/integrability data & \status{S} \\
Jacobian & Torus of line bundles & Abelianized phase/integrable variables & \status{S} \\
Picard group & Line-bundle classes & Phase sectors or charge classes & \status{S/H} \\
Hodge structure & Cohomological decomposition & Analytic-topological complexity & \status{S} \\
Mixed Hodge structure & Layered singular/relative cohomology & Amplitudes with boundaries/degenerations & \status{S} \\
Variation of Hodge str. & Periods varying over parameters & Amplitudes over kinematic space & \status{S} \\
Gauss--Manin connection & Transport of cohomology over moduli & Differential equations for amplitudes & \status{S} \\
Picard--Fuchs equation & Period differential equation & Equation satisfied by integral/amplitude & \status{S} \\
Riemann--Hilbert corr. & ODEs to monodromy & Dynamics to analytic-continuation data & \status{S} \\
Monodromy & Transport around singularity & Branch cuts or analytic continuation & \status{S} \\
Discriminant locus & Degenerate parameter values & Thresholds or singular kinematics & \status{S} \\
Positive geometry & Region with canonical form & Geometry encoding scattering amplitude & \status{S} \\
Canonical form & Distinguished differential form & Amplitude or integrand form & \status{S} \\
Amplituhedron & Positive geometry for amplitudes & Scattering encoded geometrically & \status{S} \\
Positive Grassmannian & Positive cell geometry & On-shell amplitude combinatorics & \status{S} \\
Cluster algebra & Mutating coordinate system & Different amplitude charts/channels & \status{S/H} \\
Tropical geometry & Degeneration or asymptotic geometry & Classical, graph, or boundary regime & \status{S/H} \\
Toric variety & Combinatorial algebraic geometry & Phase space from polytope data & \status{S} \\
Newton polytope & Monomial support geometry & Semiclassical/dominant-balance structure & \status{H} \\
\end{longtable}
\end{center}

\begin{remark}[Reading the dictionary through the status monoid]
The dictionary is not a flat list: it is stratified by \Cref{def:statusmonoid}. The purely
mathematical content of every row is \status{S} (these are standard objects of algebraic geometry);
the composite labels record that the \emph{physical} reading may drop to \status{H}. For instance,
``ideal $=$ constraint equations'' is \status{S/H} because the algebra is standard but the claim that
\emph{every} physical constraint is polynomial is heuristic. When we compose dictionary entries along
the realization pipeline, \Cref{def:statusmonoid} tells us the composite status; e.g.\ realizing an
amplitude via a variation of Hodge structure (\status{S}) over a moduli \emph{space} (\status{S}) is
\status{S}, but reading the resulting object as ``the'' cosmic-Galois-covariant amplitude (a Part~II
\status{H/P} entry) forces the composite to \status{H/P}.
\end{remark}

% ============================================================
\section{Moduli, Stacks, and Gauge-Correct Possibility}
\label{sec:moduli}

\subsection{Moduli functors and stacks}

The space of ``physically distinct configurations'' of a given type is a \emph{moduli space}. We phrase
it functorially, which is both the mathematically robust formulation and the one matching our
functor-of-points encoding (\Cref{rem:fop}).

\begin{definition}[Moduli functor, moduli stack]
\label{def:moduli}
A \emph{moduli functor} for a class of objects with a notion of family is a functor
$\mathcal M\colon \Sch^{\op}\to\Set$ sending a test scheme $T$ to the set of families over $T$ modulo
isomorphism. If instead one records the isomorphisms (does not quotient by them), one obtains a
\emph{moduli stack} $\underline{\mathcal M}\colon \Sch^{\op}\to\Grpd$, valued in groupoids. The moduli
functor is \emph{represented} by a scheme (resp.\ the stack by an algebraic stack) when it satisfies
effective descent and admits a universal family.
\end{definition}

\begin{definition}[Quotient stack]
\label{def:quotstack}
Let an algebraic group $G$ act on a scheme $X$. The \emph{quotient stack} $[X/G]$ is the stackification
of the action groupoid $X\times G \rightrightarrows X$ (source and target maps $(x,g)\mapsto x$ and
$(x,g)\mapsto g\cdot x$). Its $T$-points form the groupoid of principal $G$-bundles $P\to T$ together
with a $G$-equivariant map $P\to X$. The \emph{coarse space} $X/G$ is the scheme (or algebraic space)
that best approximates the stack: it is characterized by the universal property of co-representing
$[X/G]$, i.e.\ it receives a canonical map $[X/G]\to X/G$ that is universal among maps from $[X/G]$ to
schemes. For a reductive $G$ acting on an affine $X=\Spec A$ it is the GIT quotient
$X/\!/G = \Spec(A^G)$ of invariants. We stress that $X/G$ is \emph{not} in general the sheafification of
the orbit presheaf $T\mapsto X(T)/G(T)$: when stabilizers are nontrivial that sheaf need not be
representable (e.g.\ for $\mathbb G_m\curvearrowright\AAf^1$ by scaling the invariant quotient is a single
point, whereas the fppf orbit sheaf has two points).
\end{definition}

The essential invariant retained by the stack but forgotten by the coarse space is the automorphism
(stabilizer) group of each point.

\begin{proposition}[Stabilizer is the automorphism group of a point in the stack]
\label{prop:stab}
For $x\in X$, the automorphism group of the corresponding object of the groupoid $[X/G](\Spec\coeff)$
is canonically isomorphic to the stabilizer $\Stab_G(x) = \{ g\in G : g\cdot x = x\}$:
\[
  \Aut_{[X/G]}(x) \;\cong\; \Stab_G(x).
\]
\end{proposition}
\begin{proof}
By \Cref{def:quotstack}, morphisms $x\to x'$ in the action groupoid are group elements $g\in G$ with
$g\cdot x = x'$; automorphisms are therefore $g\in G$ with $g\cdot x = x$, i.e.\ exactly $\Stab_G(x)$.
Stackification is a $2$-categorical localization that preserves automorphism groups of objects
(it only adds descent data and formal inverses to already-invertible $1$-morphisms), so the
automorphism group is unchanged in $[X/G]$.
\end{proof}

\Cref{prop:stab} is the algebraic-geometry incarnation of the master formula
$\Aut_{[X/G]}(x)\simeq\Stab_G(x)$ recorded across the whole program. It is the reason a moduli stack is
``gauge-correct'': it distinguishes a configuration with no gauge automorphisms from one stabilized by a
residual gauge subgroup, even when the two have the same image in the coarse space.

\subsection{Gauge redundancy, formally}

\begin{definition}[Gauge redundancy, after Part~I/VI]
\label{def:gauge}
A \emph{gauge redundancy} on a configuration space $X$ is the action of a group $G$ whose orbits are to
be regarded as single physical states: $x$ and $g\cdot x$ describe \emph{the same} physics. The
physically correct state space is the quotient stack $[X/G]$, not merely the set of orbits $X/G$,
whenever stabilizers carry physical information (anomalies, quantization conditions, or measure
factors).
\end{definition}

\begin{remark}[Gauge redundancy is not physical symmetry]
\label{rem:gauge-not-sym}
We reiterate the program's Limitation~5: a gauge redundancy identifies \emph{different descriptions of
the same} physical content, whereas a physical symmetry maps one physical state to a \emph{genuinely
different} one. The electromagnetic potential $A\mapsto A+d\lambda$ with field strength $F=dA$
(and $d^2=0$, so $F$ is unchanged) is the paradigm: $A$ and $A+d\lambda$ are gauge-equivalent
descriptions of one physical field configuration. Algebro-geometrically this is the statement that the
physical object is a point of a \emph{stack} of connections modulo gauge, and its automorphisms are the
residual (constant) gauge transformations.
\end{remark}

% ============================================================
\section{Hodge Theory and Kinematic-Space Amplitudes}
\label{sec:hodge}

We now formalize the machinery that lets Part~II's periods \emph{vary} over a family, which is the
technical heart of \Cref{thm:gm}.

\begin{definition}[Hodge structure]
\label{def:hodge}
A (pure) \emph{Hodge structure} of weight $w$ on a finite-dimensional $\QQ$-vector space $V$ is a
decomposition $V_\CC = V\otimes_\QQ\CC = \bigoplus_{p+q=w} V^{p,q}$ with $\overline{V^{p,q}}=V^{q,p}$.
Equivalently it is a decreasing \emph{Hodge filtration} $\Fil$ with
$V_\CC = F^p\oplus\overline{F^{w-p+1}}$. A \emph{mixed} Hodge structure adds an increasing weight
filtration $W_\bullet$ with pure Hodge structures on the graded pieces (Deligne).
\end{definition}

\begin{definition}[Variation of Hodge structure]
\label{def:vhs}
A \emph{variation of Hodge structure} (VHS) over a complex base $S$ is a local system $\mathbb V$ of
$\QQ$-vector spaces on $S$ together with a holomorphic Hodge filtration $\Fil\subset \mathbb
V\otimes\Ohm_S$ such that
\begin{enumerate}
  \item each fiber $(\mathbb V_s,F^\bullet_s)$ is a Hodge structure, and
  \item \textbf{Griffiths transversality} holds: $\GM F^p \subseteq F^{p-1}\otimes\Omega^1_S$, where
  $\GM$ is the flat \emph{Gauss--Manin connection} determined by the local system.
\end{enumerate}
\end{definition}

\begin{definition}[Period map and Picard--Fuchs system]
\label{def:pf}
Let $\pi\colon\mathcal X\to S$ be a smooth projective family and $\omega_s$ a family of holomorphic
forms with locally constant integration cycles $\gamma(s)\in H_\bullet(\mathcal X_s)$. The
\emph{period} is $\Pi(s)=\int_{\gamma(s)}\omega_s$. Because the cycles $\gamma(s)$ are flat for $\GM$
(they are the flat sections) while $[\omega_s]$ is not, the period $\Pi$ is a non-constant function
satisfying a linear ODE system --- the \emph{Picard--Fuchs system} --- obtained by expressing the
high-order Gauss--Manin derivatives $[\omega_s],\GMd[\omega_s],\dots,(\GMd)^r[\omega_s]$ in the
(finite-rank, rank-$r$) de Rham cohomology basis to produce a linear dependence and pairing the result
with $\gamma(s)$:
\[
  \sum_{j=0}^{r} c_j(s)\,(\GMd)^j[\omega_s] = 0
  \quad\Longrightarrow\quad
  \mathcal L_{\mathrm{PF}}\,\Pi = 0,\ \ \mathcal L_{\mathrm{PF}}=\sum_{j=0}^r c_j(s)\,\partial^j,
\]
a differential operator of order $\le \rk H^\bullet_{\mathrm{dR}}(\mathcal X_s)$. (Note $\GM\Pi=0$ does
\emph{not} hold: the period is not a flat section; see \Cref{thm:gm}.)
\end{definition}

\begin{example}[Legendre family of elliptic curves, \status{S}]
\label{ex:legendre}
The family $E_\lambda\colon y^2 = x(x-1)(x-\lambda)$ over $S=\PP^1\setminus\{0,1,\infty\}$ has periods
$\Pi(\lambda)=\oint \frac{dx}{y}$ satisfying the hypergeometric Picard--Fuchs equation
\[
  \lambda(1-\lambda)\,\Pi'' + (1-2\lambda)\,\Pi' - \tfrac14\,\Pi = 0,
\]
whose solutions are ${}_2F_1(\tfrac12,\tfrac12;1;\lambda)$ and its analytic continuation. The
discriminant locus $\{0,1,\infty\}$ is exactly where $E_\lambda$ degenerates (a node appears); the
monodromy of $\Pi$ around these points is the physical branch-cut/threshold data. This is the simplest
nontrivial instance of ``amplitudes over kinematic space'': the elliptic Feynman-integral sector.
\end{example}

% ============================================================
\section{Positive Geometries and Geometric Factorization}
\label{sec:posgeom}

\begin{definition}[Positive geometry, after Arkani-Hamed--Bai--Lam~\cite{abl}]
\label{def:posgeom}
A \emph{positive geometry} is a pair $(X,X_{\ge 0})$ of a complex variety $X$ and a real
top-dimensional region $X_{\ge 0}\subseteq X(\RR)$ with boundary components $\partial_i X_{\ge 0}$,
together with a unique \emph{canonical form} $\Omega(X_{\ge 0})$ --- a meromorphic top form with simple
poles exactly along the boundary --- satisfying the recursive \emph{residue axiom}: for each boundary
component,
\[
  \Res_{\partial_i}\,\Omega(X_{\ge 0}) \;=\; \Omega(\partial_i X_{\ge 0}),
\]
i.e.\ the residue along a facet is the canonical form of that facet, itself a positive geometry of one
lower dimension.
\end{definition}

\begin{example}[Interval and simplex, \status{S}]
For the interval $[a,b]\subset\PP^1$ the canonical form is
$\Omega = \bigl(\frac{1}{x-a}-\frac{1}{x-b}\bigr)dx = \frac{(b-a)\,dx}{(x-a)(x-b)}$, with
$\Res_{x=a}\Omega = +1 = \Omega(\{a\})$ and $\Res_{x=b}\Omega = -1$. Iterating, the standard simplex
$\Delta^n$ has canonical form the standard volume form with logarithmic singularities on its facets;
each residue is the canonical form of a facet simplex. The recursion terminates at vertices (points),
whose canonical form is the number $\pm 1$.
\end{example}

\begin{remark}[Physical reading, \status{S}]
In scattering-amplitude applications (the amplituhedron for planar $\mathcal N=4$ super Yang--Mills,
and its cousins) the canonical form \emph{is} the amplitude integrand, and the residue axiom
\emph{is} the physical factorization of the amplitude on a singularity: the residue on a boundary
component reproduces the product of lower-point amplitudes on the corresponding factorization channel.
The recursive tree of residues (\Cref{def:restree} below) is thus a geometric model of the full
singularity structure of the amplitude.
\end{remark}

\begin{definition}[Residue tree]
\label{def:restree}
For a positive geometry $(X,X_{\ge 0})$ the \emph{residue tree} $\mathcal T(X_{\ge 0})$ is the rooted
tree whose root is decorated by $\Omega(X_{\ge 0})$, whose children are the residue trees of the
boundary components $\partial_i X_{\ge 0}$, and whose leaves are the $0$-dimensional strata decorated
by $\pm 1$.
\end{definition}

The residue tree is the object we compare, in \Cref{thm:posgeom}, with the coaction tree of Part~II's
motivic amplitude objects.

% ============================================================
\section{Main Results}
\label{sec:theorems}

We now state and prove the four theorems announced in \Cref{sec:intro}. Each extends a Part~II
construction over a geometric family and each carries an explicit status label.

\subsection{T1: Gauss--Manin family realization}

\begin{theorem}[Gauss--Manin transport implements the family realization pipeline; status \status{S}]
\label{thm:gm}
Let $\pi\colon\mathcal X\to S$ be a smooth projective family of complex varieties over a smooth connected
base $S$ (so $\mathcal X$, $S$ are varieties over $\CC$, or complex manifolds), and let
$\mathcal H_\alpha = R^k\pi^{\mathrm{an}}_\ast\underline{\QQ}$ be the associated Betti local system on the
analytification $S^{\mathrm{an}}$ (Betti cohomology is a sheaf in the analytic, not the Zariski,
topology), with its Hodge-theoretic realization the variation of Hodge structure of \Cref{def:vhs}
carrying the Gauss--Manin connection $\GM$. (The restriction to the complex/analytic setting is needed for
the Betti and Hodge realizations and for Ehresmann's theorem below.) Then:
\begin{enumerate}
  \item The assignment $s\mapsto \Real_{\mathrm{Hodge}}\bigl(\Mot(\mathcal X_s)\bigr)$ is the fiber of
  $\mathcal H_\alpha$ at $s$. The Betti cycles $\gamma(s)\in H_k(\mathcal X_s;\QQ)$ are the \emph{flat}
  (locally constant) sections of the \emph{homology} local system
  $R_k\pi^{\mathrm{an}}_\ast\underline{\QQ}$, dual to the de Rham cohomology bundle
  $\mathcal H_{\mathrm{dR}}$ carrying $\GM$, while the holomorphic form
  $[\omega_s]\in F^k$ is in general \emph{not} flat: Griffiths transversality gives
  $\GM[\omega_s]\in F^{k-1}\otimes\Omega^1_S$, which is nonzero in general. The period
  $\Pi(s)=\langle\gamma(s),[\omega_s]\rangle=\int_{\gamma(s)}\omega_s$ is the pairing of the flat cycle
  with the non-flat form, and its variation is
  \[
    d\,\Pi \;=\; \int_{\gamma(s)}\GM[\omega_s],
  \]
  so $\Pi$ is generally a non-constant function of $s$. Choosing a holomorphic frame $e_1,\dots,e_r$ of
  $\mathcal H_{\mathrm{dR}}$ with connection matrix $\GM e_i = \sum_j A_{ij}\,e_j$, the vector of periods
  $\Pi_i=\int_{\gamma}e_i$ satisfies the first-order Gauss--Manin system $d\Pi = A\,\Pi$.
  \item Consequently, eliminating the frame, the scalar period $\Pi=\int_\gamma\omega$ satisfies the
  Picard--Fuchs system $\mathcal L_{\mathrm{PF}}\Pi = 0$ (\Cref{def:pf}), a differential equation of order
  at most $\rk H^k_{\mathrm{dR}}(\mathcal X_s)$.
  \item Part~II's per-object realization pipeline \eqref{eq:motchain} is the fiber over each $s\in S$
  of a \emph{family} realization pipeline
  \[
  \begin{tikzcd}[column sep=1.5cm]
    \mathcal X \arrow[r, "\underline{\Mot}"] &
    \underline{\Mot}(\mathcal X/S) \arrow[r, "\Real_\alpha"] &
    \mathcal H_\alpha \arrow[r, "\Obs=\int_{\gamma(\cdot)}"] &
    \Pi(s),\ \mathcal L_{\mathrm{PF}}\Pi=0,
  \end{tikzcd}
  \]
  in which $\Real_\alpha$ produces the VHS $\mathcal H_\alpha$ (with its holomorphic section $[\omega_s]$)
  and $\Obs$ pairs that section against the flat Betti cycles to output the period functions $\Pi(s)$.
  Thus ``amplitudes over kinematic space'' are exactly these period functions --- solutions of the
  Picard--Fuchs system of a VHS over the moduli base $S$ --- not flat sections of the VHS.
\end{enumerate}
\end{theorem}

\begin{proof}
(1) The Gauss--Manin connection is the connection on the relative de Rham cohomology bundle
$\mathcal H_{\mathrm{dR}} = R^k\pi_\ast\Omega^\bullet_{\mathcal X/S}$ characterized by declaring the
image of the topological (Betti) local system $R^k\pi_\ast\underline{\QQ}$ to be its flat sections;
equivalently it is the unique flat connection whose horizontal sections are the locally constant
cohomology classes (Katz--Oda~\cite{katzoda}). Dually, the integration cycles
$\gamma(s)\in H_k(\mathcal X_s;\QQ)$ form a flat section of the homology local system because
$H_k(\mathcal X_s)$ is locally constant in $s$ (Ehresmann: $\pi$ is a proper submersion, hence a locally
trivial fiber bundle in the analytic topology); thus $\GM^\vee\gamma=0$. By contrast the holomorphic form
represents a section $[\omega_s]$ of the Hodge sub-bundle $F^k\subseteq\mathcal H_{\mathrm{dR}}$, and
Griffiths transversality gives $\GM[\omega_s]\in F^{k-1}\otimes\Omega^1_S$, which is nonzero in general:
$[\omega_s]$ is \emph{not} flat. Differentiating the pairing $\Pi(s)=\langle\gamma(s),[\omega_s]\rangle$ by the Leibniz rule for the flat
pairing,
\[
  d\,\Pi(s) = \underbrace{\langle\GM^\vee\gamma,[\omega_s]\rangle}_{=\,0\ \text{since}\ \GM^\vee\gamma=0}
            + \langle\gamma,\GM[\omega_s]\rangle
            = \int_{\gamma(s)}\GM[\omega_s],
\]
where the first term vanishes precisely because the Betti cycle $\gamma$ is flat ($\GM^\vee\gamma=0$).
This is generally nonzero, so $\Pi$ is a non-constant function of $s$. Choosing a holomorphic frame
$e_1,\dots,e_r$ of $\mathcal H_{\mathrm{dR}}$ with $\GM e_i=\sum_j A_{ij}e_j$, the vector of periods
$\Pi_i=\int_\gamma e_i$ obeys $d\Pi_i=\int_\gamma\GM e_i=\sum_j A_{ij}\Pi_j$, i.e.\ the first-order
Gauss--Manin system $d\Pi=A\,\Pi$. (In particular the period is emphatically \emph{not} annihilated by
$\GM$; the flat sections are the Betti cycles, against which the non-flat form is measured.)

(2) Since $\mathcal H^k_{\mathrm{dR}}(\mathcal X_s)$ is finite-dimensional of rank
$r=\rk H^k_{\mathrm{dR}}(\mathcal X_s)$, for a chosen vector field $\partial$ on $S$ the $r+1$ classes
\[
  [\omega],\ \GMd[\omega],\ \dots,\ (\GMd)^{r}[\omega]
\]
lie in an $r$-dimensional space, hence are linearly dependent over the function field of $S$:
$\sum_{j=0}^{r} c_j(s)\,(\GMd)^{j}[\omega] = 0$. Pairing with the flat cycle $\gamma(s)$ and using
$\partial^j\Pi=\int_\gamma(\GMd)^j[\omega]$ (which follows from flatness of $\gamma$ as in (1)) yields
$\sum_j c_j(s)\,\partial^j\Pi = 0$, i.e.\ $\mathcal L_{\mathrm{PF}}\Pi=0$ with
$\mathcal L_{\mathrm{PF}}=\sum_j c_j(s)\partial^j$ of order $\le r$.

(3) The functor $\underline{\Mot}$ of relative motives, the Hodge realization $\Real_\alpha$ producing
$\mathcal H_\alpha$, and the observable map $\Obs=\int_{\gamma(\cdot)}$ (pairing the holomorphic section
against flat Betti cycles) compose to the displayed family pipeline; restricting to a fiber over $s$
recovers \eqref{eq:motchain} for $X=\mathcal X_s$ by construction of the relative constructions as
fiberwise ones. Each of the three constructions is standard (Gauss--Manin/VHS theory), so by
\Cref{def:statusmonoid} the composite is status \status{S}. The identification of the resulting period
functions $\Pi(s)$ --- solutions of the Picard--Fuchs system --- with ``amplitudes over kinematic space''
is the physical reading and is the standard content of the differential-equation approach to Feynman
integrals.
\end{proof}

\begin{corollary}[Amplitudes obey their Picard--Fuchs equation]
\label{cor:amp-pf}
Any physical amplitude that arises as a period of a smooth projective family over kinematic space $S$ is
annihilated by a Picard--Fuchs operator of order at most the rank of the relevant de Rham cohomology; in
particular the analytic continuation of the amplitude around a discriminant point of $S$ is controlled by
the monodromy of $\mathcal H_\alpha$ (branch cuts $=$ monodromy, \status{S}).
\end{corollary}
\begin{proof}
Immediate from \Cref{thm:gm}(2) and the Riemann--Hilbert correspondence relating the Picard--Fuchs
$\mathcal D$-module to its monodromy representation.
\end{proof}

\Cref{thm:gm} is the precise sense in which Part~III ``fibers'' Part~II: the per-object period becomes a
period \emph{function} on the moduli base, constrained by the Picard--Fuchs system, and the
``amplitude over kinematic space'' becomes a rigorously constrained object.

\subsection{T2: Stacks are strictly finer than coarse spaces}

\begin{theorem}[Moduli stack is strictly finer than the coarse space; status \status{S}]
\label{thm:stack}
Let $G$ be an algebraic group acting on a scheme $X$ over a field $\coeff$, and suppose some point
$x\in X(\coeff)$ has nontrivial stabilizer $\Stab_G(x)\neq 1$. Then:
\begin{enumerate}
  \item \textup{(No hypotheses beyond the above.)} The canonical map $q\colon [X/G]\to X/G$ to the
  coarse space is not an isomorphism in any neighbourhood of the image of $x$: it is not fully faithful
  there, because it kills the nontrivial automorphism group $\Aut_{[X/G]}(x)\cong\Stab_G(x)$
  (\Cref{prop:stab}).
  \item \textup{(Local structure.)} Assume in addition that $G$ is \emph{smooth reductive}, that
  $\operatorname{char}\coeff=0$, and that $x$ is a smooth point contained in a $G$-invariant affine open;
  equivalently, argue after passing to formal completions. Then there is a locally closed
  $\Stab_G(x)$-invariant affine slice $W\ni x$ and an \'etale (henselian) neighbourhood $U$ of $q(x)$ for
  which
  \[
    [X/G] \times_{X/G} U \;\simeq\; \bigl[\, \Spec(\Ohm^{h}_{W,x}) \,/\, \Stab_G(x)\,\bigr],
  \]
  recording exactly the deformation-and-automorphism data that $X/G$ forgets (\'etale slice / Luna-type
  local structure). Without the reductivity/characteristic hypotheses the corresponding statement still
  holds after passing to formal completions at $x$.
\end{enumerate}
\end{theorem}

\begin{proof}
(1) The coarse space $X/G$ is by construction a \emph{scheme} (or algebraic space): every point of a
scheme has trivial automorphism group. The map $q$ sends the object $x$ of $[X/G](\Spec\coeff)$ to its
image point. On automorphism groups $q$ therefore induces the map
$\Aut_{[X/G]}(x)\cong\Stab_G(x)\to \Aut_{X/G}(q(x))=1$, which is not injective when $\Stab_G(x)\neq 1$.
A morphism of stacks that is not injective on some automorphism group is not fully faithful, hence not an
isomorphism, on any open substack containing $x$.

(2) Under the stated hypotheses ($G$ smooth reductive, $\operatorname{char}\coeff=0$, and $x$ a smooth
point in a $G$-invariant affine open), Luna's \'etale slice theorem~\cite{luna} provides a locally
closed $\Stab_G(x)$-invariant affine slice $W\ni x$ transverse to the orbit $G\cdot x$. The
multiplication map $G\times_{\Stab_G(x)} W \to X$ is \'etale onto a saturated neighbourhood of the
orbit, and passing to stack quotients gives an \'etale equivalence
$[G\times_{\Stab_G(x)} W / G]\simeq [W/\Stab_G(x)]$ over a neighbourhood of $x$; taking the local
(henselian) ring of $W$ at $x$ yields the displayed local model. This is precisely the
automorphism-plus-deformation data lost by the coarse space, whose local ring at $q(x)$ is the ring of
$\Stab_G(x)$-invariants $\bigl(\Ohm^h_{W,x}\bigr)^{\Stab_G(x)}$. When the reductivity or
characteristic-zero hypotheses fail, one replaces the \'etale slice by its formal analogue: the same
local model holds after completing at $x$, working with the $\Stab_G(x)$-action on the formal
deformation neighbourhood of $x$.
\end{proof}

\begin{corollary}[$\overline{\mathcal M}_g$ must be a stack]
\label{cor:mgbar}
The Deligne--Mumford moduli of stable genus-$g$ curves $\overline{\mathcal M}_g$ has points (e.g.\
hyperelliptic curves, or the curve with extra automorphisms) with nontrivial automorphism groups; hence
by \Cref{thm:stack} it is not faithfully represented by its coarse space. Whenever curves with
automorphisms contribute physically distinct weightings --- as in the string-theoretic path integral
over the moduli of worldsheets, where the automorphism group divides the measure --- one must integrate
over the \emph{stack}, not the coarse space.
\end{corollary}
\begin{proof}
Deligne--Mumford~\cite{delignemumford} construct $\overline{\mathcal M}_g$ as a smooth proper DM stack;
the existence of automorphism-carrying curves is classical (e.g.\ the genus-$2$ curve $y^2=x^5-x$ has
automorphism group of order $>1$). Apply \Cref{thm:stack}.
\end{proof}

\Cref{thm:stack} is the algebro-geometric member of the program's four-fold ``stacks retain gauge
information'' family: Part~I posits it informally, Part~III proves it geometrically here, Part~V will
prove a type-theoretic version (groupoid quotient vs.\ set truncation), and Part~VI will prove the
descent-theoretic version. All four are the same fact at increasing rigor.

\subsection{T3: Positive-geometry residues and the motivic coaction}

\begin{theorem}[Residue tree reproduces a coaction-like decomposition; status \status{H}]
\label{thm:posgeom}
Let $(X,X_{\ge 0})$ be a positive geometry with canonical form $\Omega$ and residue tree
$\mathcal T(X_{\ge 0})$ (\Cref{def:restree}). Then:
\begin{enumerate}
  \item (Status \status{S}.) The residue operation defines a coalgebra-like ``geometric coproduct''
  \[
    \Delta_{\mathrm{geo}}\,\Omega(X_{\ge 0}) \;=\; \sum_i \Omega(\partial_i X_{\ge 0}) \otimes
    \bigl[\text{normal data at }\partial_i\bigr],
  \]
  and iterating $\Delta_{\mathrm{geo}}$ reproduces the tree $\mathcal T(X_{\ge 0})$; this operation is
  coassociative up to the ordering of iterated boundaries (the boundary strata form a poset, giving a
  graded coalgebra structure on the free vector space on strata).
  \item (Status \status{H}, conjectural for amplituhedron-type geometries.) When $\Omega$ is the
  integrand of a Feynman/scattering amplitude whose period admits a motivic lift
  $\amp^{\mathfrak m}$ (Part~II), the tree $\mathcal T(X_{\ge 0})$ is isomorphic, as a graded rooted
  tree, to the weight-graded coaction tree of $\amp^{\mathfrak m}$ under the motivic coproduct
  $\Delta(\amp^{\mathfrak m})=\sum_i \amp^{\mathfrak m}_{i,1}\otimes\amp^{\mathfrak m}_{i,2}$; i.e.\
  geometric factorization on boundaries and motivic coaction on periods are the same decomposition.
\end{enumerate}
\end{theorem}

\begin{proof}
(1) By the residue axiom (\Cref{def:posgeom}), $\Res_{\partial_i}\Omega=\Omega(\partial_i X_{\ge 0})$,
so the map $\Delta_{\mathrm{geo}}$ is well defined on the free graded vector space $C_\bullet$ spanned by
the (canonical forms of the) boundary strata, graded by codimension. Coassociativity up to poset
ordering is the statement that the iterated residue along a codimension-two stratum is independent of the
order in which the two facets are approached whenever the strata meet transversally --- a standard
property of iterated residues of forms with logarithmic (simple-pole) singularities (Leray). Hence
$(C_\bullet,\Delta_{\mathrm{geo}})$ is a graded coalgebra and its cogenerated tree is $\mathcal T(X_{\ge
0})$. This part is a theorem of positive-geometry theory and is status \status{S}.

(2) This is the conjectural identification and we flag it \status{H}. The evidence is: (a) both sides are
weight/codimension-graded rooted trees whose leaves are $\QQ$-numbers (residues $=\pm1$ on one side,
weight-$0$ periods $=$ rationals on the other); (b) for one-loop integrands and for the polygon/associahedron
geometries the two trees are known to agree, matching the Goncharov coproduct on the associated
polylogarithms; (c) the ``cosmic Galois''/coaction principle of Brown and the residue recursion of
Arkani-Hamed--Bai--Lam are both governed by the same weight filtration. A general proof would require a
functor from the poset of boundary strata of $X_{\ge 0}$ to the comodule of motivic periods intertwining
$\Delta_{\mathrm{geo}}$ with $\Delta$; constructing this functor for all amplituhedron-type geometries is
open. We therefore assert (2) as a status-\status{H} candidate theorem, not an established one, and record
it in \Cref{sec:discussion} as the program's sharpest open problem in this domain.
\end{proof}

\begin{remark}
The honest status bookkeeping of \Cref{thm:posgeom} is exactly what the \status{S}/\status{H}/\status{P}
discipline is for: part~(1) is a genuine theorem, part~(2) is a precise conjecture. Collapsing the two
would be the kind of overclaim the program is designed to prevent (Limitation~1).
\end{remark}

\subsection{T4: The derived critical locus as BV on-shell algebra}

\begin{theorem}[Derived critical locus formalizes $R/\Ideal$ at the BV level; status \status{S}]
\label{thm:dcrit}
Let $X$ be a smooth scheme and $S\colon X\to\AAf^1$ an action functional (regular function). Let
$\Crit(S)=\Spec\bigl(\Ohm_X/(dS)\bigr)$ be the classical critical locus and $\dCrit(S)$ the
\emph{derived} critical locus, i.e.\ the derived zero locus of $dS\in\Gamma(X,\Omega^1_X)$, constructed
as the derived intersection of the graph of $dS$ with the zero section inside $T^\ast X$:
\[
  \dCrit(S) \;=\; X \times^{h}_{T^\ast X}\, X.
\]
Then:
\begin{enumerate}
  \item The classical truncation is $\pi_0(\dCrit(S)) = \Crit(S)$; in particular
  $H^0(\Ohm_{\dCrit(S)}) = \Ohm_X/(dS) = R/\Ideal$ with $\Ideal=(dS)$ the ideal of equations of motion.
  Thus the on-shell observable algebra of \Cref{ex:onshell} is the $H^0$ of the derived object.
  \item $\dCrit(S)$ carries a canonical $(-1)$-shifted symplectic structure (Pantev--To\"en--Vaqui\'e--Vezzosi),
  and its structure sheaf, computed via the Koszul resolution of $(dS)$, is the field/antifield sector of
  the Batalin--Vilkovisky complex: the shifted symplectic form is the BV antibracket, the degree-$0$
  generators are the fields, and the degree-$(-1)$ (and lower) generators are the antifields. No ghosts
  (positive cohomological degree) appear, because $X$ is a scheme with no gauge symmetry; ghosts are
  produced only when $X$ is upgraded to an algebraic \emph{stack} $[X/G]$ resolving a gauge redundancy
  (see \Cref{rem:ghosts}).
\end{enumerate}
Consequently the dictionary entry ``derived critical locus $\leftrightarrow$ perturbative field theory
around the action'' is upgraded from \status{S/H} to \status{S} for the field/antifield sector.
\end{theorem}

\begin{proof}
(1) The section $dS$ of $\Omega^1_X$ has classical zero locus $\Crit(S)=V(dS)$; its coordinate ring is
$\Ohm_X/(dS)$, which is exactly the on-shell algebra $R/\Ideal$ of \Cref{ex:onshell} for
$\Ideal=(\partial_i S)$. Derived intersection only refines the vanishing locus by remembering the Koszul
syzygies; passing to $\pi_0$ (equivalently $H^0$ of the structure complex) discards the higher Koszul
terms and returns the classical quotient. Hence $\pi_0(\dCrit(S))=\Crit(S)$ and $H^0=R/\Ideal$.

(2) Model $\dCrit(S)$ by the Koszul complex $\bigl(\Ohm_X\otimes \Lambda^\bullet T_X,\, \iota_{dS}\bigr)$,
the Koszul complex contracting with $dS$; this is the algebra of polyvector fields with differential
$\iota_{dS}=\{S,-\}$, i.e.\ the field/antifield part of the BV complex, with $\Lambda^p T_X$ concentrated
in cohomological degree $-p$: degree $0$ carries the fields $\Ohm_X$ and degrees $-1,-2,\dots$ carry the
antifields $T_X,\Lambda^2 T_X,\dots$. Every generator sits in \emph{nonpositive} degree; there are no
positive-degree ghosts, consistent with the absence of gauge symmetry for a bare scheme $X$. The
$(-1)$-shifted cotangent identification $\dCrit(S)=T^\ast[-1]X|_{dS}$ endows it with the canonical
$(-1)$-shifted symplectic form of PTVV~\cite{ptvv}; unwinding the pairing on polyvector fields gives the
odd Poisson bracket, which is the BV antibracket. All ingredients (Koszul resolution, PTVV shifted
symplectic structures, the identification of the field/antifield BV complex with functions on the
$(-1)$-shifted cotangent of the critical locus) are established mathematics; see also Costello--Gwilliam's
factorization-algebra formulation~\cite{costellogwilliam}. Hence the upgrade to status \status{S}.
\end{proof}

\begin{remark}[Where ghosts come from: schemes give antifields, stacks give ghosts]
\label{rem:ghosts}
It is important not to overstate \Cref{thm:dcrit}. For a bare smooth \emph{scheme} $X$ the Koszul complex
$\Lambda^\bullet T_X$ lives entirely in nonpositive cohomological degrees, so the derived critical locus
supplies only the \emph{field/antifield} sector (degrees $0$ and $\le -1$) of the BV complex --- there is
no gauge symmetry and hence no ghosts. Ghosts (positive cohomological degree) arise precisely when one
resolves a \emph{gauge redundancy}: replacing $X$ by an algebraic stack $[X/G]$ introduces, via the
Chevalley--Eilenberg differential of the Lie algebra of $G$, the positive-degree ghost generators
$\mathfrak g[1]$ dual to the gauge directions. Thus, in a slogan, \emph{derived geometry supplies the
antifields, stacky geometry supplies the ghosts}; the full BV--BRST complex of a gauge theory is the
structure sheaf of the derived critical locus \emph{of the stack} $[X/G]$. This is exactly the point at
which \Cref{thm:stack} (stacks retain gauge data) and \Cref{thm:dcrit} (derived critical loci) meet.
\end{remark}

\begin{remark}[Why the upgrade matters]
The primary source lists ``derived critical locus $\leftrightarrow$ BV/BRST field space'' as \status{S/H}
because, as a bare dictionary slogan, the physical reading was heuristic. \Cref{thm:dcrit} shows the
field/antifield part of the slogan is in fact a theorem of derived algebraic geometry: the field/antifield
BV complex is \emph{literally} the structure sheaf of the derived critical locus, and (by
\Cref{rem:ghosts}) the ghost sector is recovered by the same construction applied to the gauge stack. This
is a model of what the whole program aspires to --- converting heuristic dictionary entries into theorems
where the mathematics permits --- and it is the cleanest such upgrade in the algebraic-geometry domain.
\end{remark}

% ============================================================
\section{The Family Realization Pipeline}
\label{sec:family}

We assemble the theorems into the promised generalization of Part~I's pipeline from a single object to a
moduli-indexed family, and we track statuses with \Cref{def:statusmonoid}.

\begin{definition}[Family realization pipeline]
\label{def:familypipeline}
A \emph{family realization pipeline} over a base scheme $S$ is the data
\[
  \bigl(\pi\colon\mathcal X\to S,\ \underline{\Mot},\ \Real_\alpha,\ \GM,\ \Obs\bigr)
\]
such that $\Real_\alpha\circ\underline{\Mot}$ produces a VHS $\mathcal H_\alpha$ on $S$ with Gauss--Manin
connection $\GM$ and holomorphic section $[\omega_s]$, and the observable
$\Obs=\int_{\gamma(\cdot)}$ pairs that section against the flat Betti cycles $\gamma(s)$. The output is a
map
\[
  \Obs_\alpha^{S}\colon\ \{\text{families over }S\}\ \longrightarrow\ \mathrm{Sol}(\mathcal L_{\mathrm{PF}}),
  \qquad \Obs_\alpha^S(\mathcal X) = \Pi(s)\ \text{with}\ \mathcal L_{\mathrm{PF}}\Pi=0,
\]
where $\mathrm{Sol}(\mathcal L_{\mathrm{PF}})$ is the (finite-dimensional) local solution space of the
Picard--Fuchs system of $\mathcal H_\alpha$. (The dual statement: integrating a \emph{holomorphic
cohomology frame} $e_1,\dots,e_r$ against a \emph{single} flat cycle $\gamma$ yields a period vector that
is a horizontal section of the connection dual to $\GM$; integrating one non-flat form against a flat
homology frame yields the non-horizontal, Picard--Fuchs-constrained period functions above.)
\end{definition}

\begin{proposition}[The family pipeline reduces fiberwise to Part~I's pipeline; status \status{S}]
\label{prop:reduce}
For each $s\in S$, evaluating the period function $\Obs_\alpha^S(\mathcal X)=\Pi$ at $s$ recovers Part~I's
$\Obs_\alpha(\mathcal X_s) = \Obs(\Real_\alpha(\Phi(\mathcal X_s)))$ with $\Phi=\Mot$. Hence the family
pipeline is a genuine extension: it agrees with Part~I/II fiberwise and adds globally the Picard--Fuchs
constraint $\mathcal L_{\mathrm{PF}}\Pi=0$ (equivalently the first-order Gauss--Manin system on the period
vector).
\end{proposition}
\begin{proof}
Immediate from \Cref{thm:gm}(2)--(3): the relative constructions $\underline{\Mot}$, $\Real_\alpha$,
$\Obs$ are fiberwise the absolute ones, and evaluation
$\mathrm{Sol}(\mathcal L_{\mathrm{PF}})\to \CC$, $\Pi\mapsto\Pi(s)$, gives back the single-object
observable at $s$.
\end{proof}

\begin{proposition}[Status of the family pipeline]
\label{prop:familystatus}
The family realization pipeline is status \status{S} whenever $\pi$ is smooth projective and the
realization channel $\alpha$ is one of $\{$Betti, de Rham, Hodge$\}$; it drops to \status{S/H} for the
\'etale channel over a non-closed base field (arithmetic monodromy is only heuristically an
``observable''), and to \status{H/P} if one further posits a cosmic-Galois action on the period-solution
space $\mathrm{Sol}(\mathcal L_{\mathrm{PF}})$ (a Part~II \status{H/P} entry).
\end{proposition}
\begin{proof}
By \Cref{def:statusmonoid} the status of the composite is the $\wedge$-join of the component statuses.
Gauss--Manin/VHS theory over $\CC$ is \status{S}; the \'etale realization's identification with a
physical observable is \status{S/H} (arithmetic, not directly measured); the cosmic-Galois covariance is
\status{H/P} (Part~II, T5). Joining gives the stated results.
\end{proof}

\Cref{prop:familystatus} is exactly the epistemic bookkeeping the program demands: it names precisely
which realization channels keep the pipeline standard and which push it into heuristic or speculative
territory.

% ============================================================
\section{Worked Examples}
\label{sec:examples}

\subsection{The Legendre family as a family pipeline}

Combining \Cref{ex:legendre} with \Cref{def:familypipeline}: for $\pi\colon\mathcal E\to
S=\PP^1\setminus\{0,1,\infty\}$ the Legendre family, the VHS $\mathcal H = R^1\pi_\ast\underline{\QQ}$
has rank $2$, the periods $\Pi(\lambda)$ satisfy the hypergeometric Picard--Fuchs operator, and the
monodromy representation $\pi_1(S)\to SL_2(\ZZ)$ around $\{0,1,\infty\}$ is the branch-cut data. This is a
complete status-\status{S} instance of the entire Part~III machinery: a moduli base, a VHS, a
Picard--Fuchs equation, and monodromy as physical analytic continuation.

\subsection{\texorpdfstring{A quotient stack: $[\AAf^1/\mathbb G_m]$}{A quotient stack of the affine line by the multiplicative group}}

Let $\mathbb G_m$ act on $\AAf^1$ by scaling. The GIT/coarse quotient is
$\Spec\bigl(\coeff[x]^{\mathbb G_m}\bigr)=\Spec\coeff$, a \emph{single} point (the only invariant
functions are the constants, and the only closed orbit is the origin), so the coarse space records
nothing beyond a point. The stack $[\AAf^1/\mathbb G_m]$, by contrast, has two points: the open orbit,
whose objects have trivial automorphisms, and the origin, whose automorphism group is
$\Aut(0)\cong\mathbb G_m$ (by \Cref{prop:stab}, $\Stab_{\mathbb G_m}(0)=\mathbb G_m$). This is the
algebraic model of a $U(1)$ gauge theory at a symmetric point: the residual gauge group at the origin is
retained by the stack and forgotten by the coarse space, exactly as \Cref{thm:stack} predicts.

\subsection{A derived critical locus: the free scalar}

For $X=\AAf^1$ with coordinate $\phi$ and action $S(\phi)=\tfrac{1}{2}m^2\phi^2$, we have $dS=m^2\phi\,
d\phi$, $\Crit(S)=\{ \phi=0\}$, and $R/\Ideal = \coeff[\phi]/(\phi)=\coeff$: the unique on-shell state is
$\phi=0$. The derived critical locus resolves the (here already reduced) intersection by the Koszul
complex $\coeff[\phi]\xleftarrow{\ \cdot m^2\phi\ }\coeff[\phi]\,\xi$ with $\xi$ the \emph{antifield} in
degree $-1$; the $(-1)$-shifted symplectic form pairs $\phi$ with $\xi$, giving the BV antibracket
$\{\phi,\xi\}=1$. There are no ghosts, since the free scalar has no gauge symmetry (consistent with
\Cref{rem:ghosts}). This is the smallest nontrivial illustration of \Cref{thm:dcrit}.

% ============================================================
\section{Formalization: Haskell and Lean}
\label{sec:code}

We accompany the paper with executable \textsf{Haskell} encodings of the four load-bearing structures
(scheme as functor of points, quotient stack, variation of Hodge structure, positive geometry) and a
best-effort \textsf{Lean~4} sketch of the same type signatures. The Haskell modules compile with GHC and
their \texttt{main} prints demonstrations of \Cref{prop:stab,thm:posgeom}(1), and the residue-tree
construction. Full sources are in
\texttt{src/\allowbreak algebraic-\allowbreak geometry-\allowbreak physical-\allowbreak possibility/}
and
\texttt{lean/\allowbreak algebraic-\allowbreak geometry-\allowbreak physical-\allowbreak possibility/}.

\subsection{Scheme as functor of points}

Following \Cref{rem:fop} we model a scheme by its functor of points $\CRing\to\Set$, and the on-shell
algebra $R/\Ideal$ by a coordinate ring with a distinguished constraint ideal:
\begin{verbatim}
-- A commutative ring presented by generators and a constraint ideal.
data CoordRing = CoordRing { generators :: [String], relations :: [Poly] }

-- A scheme exposed via its functor of points R |-> Hom(Spec R, X).
newtype Scheme = Scheme { pointsOver :: CoordRing -> [RPoint] }

-- The on-shell algebra R/I: functions modulo the ideal of constraints.
-- (idealGens is the record accessor for an Ideal's generating polynomials.)
onShell :: CoordRing -> Ideal -> CoordRing
onShell r i = r { relations = relations r ++ idealGens i }
\end{verbatim}

\subsection[Quotient stack with stabilizer data]{Quotient stack with stabilizer data (\Cref{prop:stab})}

\begin{verbatim}
-- A finite group action, enough to compute stabilizers concretely.
data GAction g x = GAction { act :: g -> x -> x, elements :: [g] }

-- Stabilizer of a point: {g | g . x == x}. This is Aut_{[X/G]}(x).
stabilizer :: (Eq x) => GAction g x -> x -> [g]
stabilizer ga x = [ g | g <- elements ga, act ga g x == x ]

-- A quotient stack remembers, for each point, its automorphism group.
data QuotientStack g x = QuotientStack
  { action    :: GAction g x
  , autGroup  :: x -> [g]        -- x |-> Stab_G(x) = Aut_{[X/G]}(x)
  }
\end{verbatim}

The demonstration in \texttt{Main.hs} builds $[\ZZ/4 \curvearrowright \ZZ/4]$ (rotation of a square) and
prints the automorphism group at the fixed point, exhibiting \Cref{prop:stab} concretely: the coarse
quotient forgets it, the stack does not.

\subsection{Variation of Hodge structure with Gauss--Manin connection}

\begin{verbatim}
-- A VHS over a base: a fiber vector space, a Hodge filtration, and a
-- flat (Gauss-Manin) connection whose FLAT SECTIONS are the Betti cycles
-- (a basis against which the non-flat holomorphic form is measured).
data VHS base = VHS
  { fiber        :: base -> [Double]      -- local system fiber
  , hodgeFiltr   :: base -> [[Double]]    -- F^p, decreasing filtration
  , gaussManin   :: Connection base       -- flat connection nabla
  }

-- Derive the Picard-Fuchs operator by expressing the iterated
-- Gauss-Manin derivatives (nabla^j [omega]) in the finite-rank
-- de Rham basis and reading off the linear dependence. It
-- ANNIHILATES the period function Pi(s) -- the pairing of a flat
-- cycle with the NON-flat form -- which is not a flat section.
derivePicardFuchs :: VHS base -> ODE
derivePicardFuchs = linDep . gmDerivatives
\end{verbatim}

\subsection{Positive geometry and the recursive residue tree}

\begin{verbatim}
-- A positive geometry: a region, its canonical form, and its boundary
-- strata (themselves positive geometries) -- the recursion of Def 7.1.
data PositiveGeometry = PositiveGeometry
  { canonicalForm :: DiffForm
  , boundaries    :: [PositiveGeometry]
  }

-- The residue tree of Def 7.4 / Theorem 8.5(1): root = canonical form,
-- children = residue trees of the boundary strata.
data Tree a = Node a [Tree a]

residueTree :: PositiveGeometry -> Tree DiffForm
residueTree pg = Node (canonicalForm pg)
                      (map residueTree (boundaries pg))
\end{verbatim}

\subsection{Lean sketch}

The \textsf{Lean~4} file records idiomatic type signatures for the same structures (scheme via a functor
of points, quotient stack with a stabilizer field, VHS with Griffiths transversality, positive geometry
with the residue axiom) and states \Cref{thm:stack} as a signature with a \texttt{sorry} placeholder. It
is a best-effort sketch, not build-gated:
\begin{verbatim}
structure QuotientStack (X G : Type) [Group G] [MulAction G X] where
  autGroup : X -> Subgroup G            -- x |-> Stab_G(x) = Aut_{[X/G]}(x)
  coarseMap : X -> Quotient (MulAction.orbitRel G X)

-- T2: a nontrivial stabilizer forces the coarse map to lose information.
theorem stack_finer_than_coarse
    {X G : Type} [Group G] [MulAction G X] (x : X)
    (h : MulAction.stabilizer G x /= (bot : Subgroup G)) :
    True := by trivial   -- best-effort placeholder; see paper Thm 8.3
\end{verbatim}

% ============================================================
\section{Discussion}
\label{sec:discussion}

\subsection{How Part~III respects the standing limitations}
\label{sec:limitations}

The program carries five standing limitations from Part~I; we discharge each for the
algebraic-geometry domain.

\begin{enumerate}
  \item \emph{Not every analogy is a physical theory.} We enforce this with the status labels of
  \Cref{tab:dict1,tab:dict2} and, sharply, with \Cref{thm:posgeom}: its part~(1) is \status{S} and its
  part~(2) is \status{H}, and we refuse to collapse them.
  \item \emph{Motives are not literally ``the functions of'' an object.} We used $\Mot$ only through its
  \emph{realizations} $\Real_\alpha$ (Hodge, Betti, de Rham), which are standard functors; the
  ``functional essence'' reading remains a Part~II \status{H} slogan and is never invoked as a premise in
  a proof here.
  \item \emph{No single universal functor $\mathrm{Math}\to\mathrm{Phys}$.} Our pipeline is
  domain-local: it is a functor \emph{on the algebraic-geometry domain} $\Math_{\mathsf{AG}}$, fibered
  over a base $S$, not a global functor. \Cref{def:familypipeline} is explicitly indexed by a family and
  a channel $\alpha$.
  \item \emph{Not every observable is a period; elliptic/modular/non-Tate structures intervene.}
  \Cref{ex:legendre} already exhibits an \emph{elliptic} period, and \Cref{thm:posgeom}(2) is flagged
  \status{H} precisely because non-Tate geometries obstruct a naive polylogarithmic coaction (Part~II,
  T4). Corollary~\ref{cor:amp-pf} bounds only the \emph{order} of the Picard--Fuchs operator, not the
  transcendental type of its solutions.
  \item \emph{Gauge redundancy $\neq$ physical symmetry.} \Cref{rem:gauge-not-sym} states this
  explicitly; \Cref{thm:stack} makes it geometric (stabilizers are gauge automorphisms, retained by the
  stack, not new physical states).
\end{enumerate}

\subsection{Open problems handed forward}
\label{sec:open}

\begin{itemize}
  \item \textbf{Positive geometry $=$ motivic coaction (\Cref{thm:posgeom}(2)).} The sharpest open
  problem of this domain: construct, for all amplituhedron-type geometries, a functor from the poset of
  boundary strata to the comodule of motivic periods intertwining $\Delta_{\mathrm{geo}}$ with the
  motivic coproduct $\Delta$. Partial results exist for one-loop integrands and polygon/associahedron
  geometries; the general case is open (``Hopf algebra of the amplituhedron'').
  \item \textbf{Stacky path integrals.} Making \Cref{cor:mgbar} quantitative (the precise measure on
  $\overline{\mathcal M}_g$ weighted by $1/|\!\Aut|$) connects Part~III to Part~IV's TQFT functors and
  Part~VI's stacks-and-descent treatment.
  \item \textbf{Derived enhancement of the whole dictionary.} \Cref{thm:dcrit} upgrades one row to
  \status{S} via derived geometry; systematically deriving the tangent/cotangent-complex rows of
  \Cref{tab:dict2} (BV/BRST for gauge theories) is a natural next step toward Part~VI.
\end{itemize}

\subsection{How Part~III seeds Part~IV}

Part~IV (algebraic topology: conserved global information) takes the varieties and moduli produced here
and studies their \emph{topological invariants} stable under continuous deformation: the (co)homology
$H_\bullet(X)$ underlying the Betti realization of \Cref{thm:gm}, the characteristic classes of the
bundles in \Cref{tab:dict1}, and the bordism/TQFT functors that make realization functorial. In
particular, the monodromy representation of \Cref{cor:amp-pf} is a topological (local-system) datum, and
the automorphism data of \Cref{thm:stack} is what Part~IV's groupoid-cardinality/TQFT weightings act on.
Part~III thus supplies the geometric substrate; Part~IV supplies its topology.

% ============================================================
\section{Conclusion}
\label{sec:conclusion}

We have formalized the algebraic-geometry layer of the Math$\to$Physics Representation Library as the
\emph{geometry of physical possibility}, preserving the program's \status{S}/\status{H}/\status{P}
status discipline throughout. The central technical contribution is the extension of Part~II's
per-object realization pipeline to a \emph{family} pipeline fibered over a moduli base: Gauss--Manin
transport realizes periods as the pairing of flat Betti cycles with the non-flat holomorphic form, giving
period functions constrained by Picard--Fuchs equations (\Cref{thm:gm}), moduli stacks retain the gauge
automorphisms that coarse spaces forget (\Cref{thm:stack}), positive
geometries encode factorization through recursive residues that conjecturally match the motivic coaction
(\Cref{thm:posgeom}), and derived critical loci realize the on-shell algebra $R/\Ideal$ at the BV level,
upgrading a heuristic dictionary entry to a theorem (\Cref{thm:dcrit}).

In the program's compressed slogans:
\begin{quote}
mathematics is the syntax of physical representation;\\
motives are functional semantics;\\
periods are numerical measurements;\\
coalgebras are decomposition laws.
\end{quote}
Part~III adds a fifth clause: \emph{geometry is the space in which those measurements vary}. The moduli
space is the stage, its singularities are the thresholds, its compactification supplies the asymptotic
sectors, and its stackiness is the gauge redundancy: a single algebro-geometric object carrying, in
its points, its automorphisms, and its variation of Hodge structure, the entire apparatus of physical
possibility. Part~IV will endow this stage with its conserved topological invariants.

% ============================================================
\begin{thebibliography}{99}

\bibitem{primary}
Anonymous Draft,
\emph{A Representation-Stack Library for Mathematical Structures in Physical Representation: Motives,
Periods, Coalgebras, Gauge Redundancy, and Quantum-Informational Availability}.
Primary source of the Math$\to$Physics Representation Library program, 2024.

\bibitem{hartshorne}
R.~Hartshorne,
\emph{Algebraic Geometry}.
Graduate Texts in Mathematics \textbf{52}, Springer-Verlag, New York, 1977.

\bibitem{mumfordgit}
D.~Mumford, J.~Fogarty, and F.~Kirwan,
\emph{Geometric Invariant Theory}, 3rd ed.
Ergebnisse der Mathematik und ihrer Grenzgebiete, Springer, Berlin, 1994.

\bibitem{delignemumford}
P.~Deligne and D.~Mumford,
``The irreducibility of the space of curves of given genus,''
\emph{Publications Math\'ematiques de l'IH\'ES} \textbf{36} (1969), 75--109.

\bibitem{vistoli}
A.~Vistoli,
``Notes on Grothendieck topologies, fibered categories and descent theory,''
arXiv:math/0412512 (2004).

\bibitem{abl}
N.~Arkani-Hamed, Y.~Bai, and T.~Lam,
``Positive geometries and canonical forms,''
\emph{Journal of High Energy Physics} \textbf{11} (2017) 039; arXiv:1703.04541.

\bibitem{amplituhedron}
N.~Arkani-Hamed and J.~Trnka,
``The Amplituhedron,''
\emph{Journal of High Energy Physics} \textbf{10} (2014) 030; arXiv:1312.2007.

\bibitem{costellogwilliam}
K.~Costello and O.~Gwilliam,
\emph{Factorization Algebras in Quantum Field Theory}, Vol.~1.
New Mathematical Monographs \textbf{31}, Cambridge University Press, 2016.

\bibitem{ptvv}
T.~Pantev, B.~To\"en, M.~Vaqui\'e, and G.~Vezzosi,
``Shifted symplectic structures,''
\emph{Publications Math\'ematiques de l'IH\'ES} \textbf{117} (2013), 271--328; arXiv:1111.3209.

\bibitem{kontsevichzagier}
M.~Kontsevich and D.~Zagier,
``Periods,''
in \emph{Mathematics Unlimited --- 2001 and Beyond}, Springer, Berlin, 2001, 771--808.

\bibitem{voisin}
C.~Voisin,
\emph{Hodge Theory and Complex Algebraic Geometry I, II}.
Cambridge Studies in Advanced Mathematics \textbf{76, 77}, Cambridge University Press, 2002--2003.

\bibitem{griffiths}
P.~Griffiths,
``Periods of integrals on algebraic manifolds, I \& II,''
\emph{American Journal of Mathematics} \textbf{90} (1968), 568--626, 805--865.

\bibitem{katzoda}
N.~Katz and T.~Oda,
``On the differentiation of de Rham cohomology classes with respect to parameters,''
\emph{Journal of Mathematics of Kyoto University} \textbf{8} (1968), 199--213.

\bibitem{luna}
D.~Luna,
``Slices \'etales,''
\emph{Bulletin de la Soci\'et\'e Math\'ematique de France, M\'emoire} \textbf{33} (1973), 81--105.

\bibitem{brown}
F.~Brown,
``Notes on motivic periods,''
\emph{Communications in Number Theory and Physics} \textbf{11} (2017), 557--655; arXiv:1512.06409.

\bibitem{goncharov}
A.~B.~Goncharov, ``Multiple polylogarithms and mixed Tate motives,'' arXiv:math/0103059 (2001);
see also \emph{Duke Mathematical Journal} \textbf{128} (2005), no.~2, 209--284.

\bibitem{part1}
M.~Long and The YonedaAI Collaboration,
\emph{Foundations: The Representation Stack and Realization Pipeline}.
Part~I of \emph{A Math$\to$Physics Representation Library}, YonedaAI Research Collective, 2026.

\bibitem{part2}
M.~Long and The YonedaAI Collaboration,
\emph{Motives, Periods, and Amplitudes: The Coalgebraic Anatomy of Observables}.
Part~II of \emph{A Math$\to$Physics Representation Library}, YonedaAI Research Collective, 2026.

\end{thebibliography}

\end{document}
